\section{Sensitivity of $\Sprj$ to $\Omega_m$, $\sigma_8$}
\label{sec:cosmo}

\textbf{Script:} \pyfile{scripts/make\_cosmo\_sensitivity.py}. The
selected/RND ratio --- the observable scored in \S\ref{sec:mock} --- at
fixed $\beff$, varying $\Omega_m$ and $\sigma_8$ independently by
$\pm5\%,\pm10\%$ around the Buzzard fiducial, at $R\in\{1,3,10,20\}\,\hMpc$
in two representative bins ($\lambda[20,30)$ and $\lambda[60,500)$,
$z\approx0.5$):

\begin{figure}[H]
\centering
\includegraphics[width=0.85\textwidth]{fig_cosmo_sensitivity.pdf}
\caption{Fractional change in the selected/RND ratio relative to the
Buzzard fiducial, vs.\ $\Omega_m$ (left) and $\sigma_8$ (right), at each
$R$ tested.}
\label{fig:cosmo}
\end{figure}

The ratio is remarkably insensitive to cosmology: across every $R$,
richness bin, and parameter tested, the maximum fractional shift under a
$\pm10\%$ cosmology change is $\CosmoMaxAbsRatioShift$. $\sigma_8$
dominates over $\Omega_m$ (more clustering amplitude strengthens the
correlated channel the boost model responds to); $\Omega_m$'s effect is
smaller and sign-mixed. Both are largest near the aperture ($R=1$) and
decay toward the two-halo regime, since the \emph{ratio} cancels most of
the direct $P(k)$ dependence between its numerator and denominator,
leaving only the residual from $\bsel(\theta)$'s own cosmology dependence
through $\beff$/$\DRND$.

This means the systematic-uncertainty budget for the selected/RND
observable is dominated by the mass-observable relation
(\S\ref{sec:alpha}) and the closure's $\delta$ estimate
(\S\ref{sec:methods}), not by cosmological parameter uncertainty ---
consistent with \S\ref{sec:methods}'s note that the Planck-vs-Buzzard
cosmology bug fixed this session was not, on its own, the source of the
residuals in \S\ref{sec:fig6}. \emph{Scope:} $\beff$ was held fixed
across the cosmology grid; letting it respond self-consistently to
$(\Omega_m,\sigma_8)$ could add a comparable secondary effect not
captured here.
